\section{Related Work}
\label{sec:related}

\noindent\textbf{Residual Representations.}
In image recognition, VLAD [18] is a representation that encodes by the residual vectors with respect to a dictionary, and Fisher Vector [30] can be formulated as a probabilistic version [18] of VLAD.
Both of them are powerful shallow representations for image retrieval and classification [4,48].
For vector quantization, encoding residual vectors [17] is shown to be more effective than encoding original vectors.

In low-level vision and computer graphics, for solving Partial Differential Equations (PDEs), the widely used Multigrid method [3] reformulates the system as subproblems at multiple scales, where each subproblem is responsible for the residual solution between a coarser and a finer scale. An alternative to Multigrid is hierarchical basis preconditioning [45,46], which relies on variables that represent residual vectors between two scales. It has been shown [3,45,46] that these solvers converge much faster than standard solvers that are unaware of the residual nature of the solutions. These methods suggest that a good reformulation or preconditioning can simplify the optimization.

\vspace{6pt}
\noindent\textbf{Shortcut Connections.}
Practices and theories that lead to shortcut connections [2,34,49] have been studied for a long time.
An early practice of training multi-layer perceptrons (MLPs) is to add a linear layer connected from the network input to the output [34,49]. In [44,24], a few intermediate layers are directly connected to auxiliary classifiers for addressing vanishing/exploding gradients. The papers of [39,38,31,47] propose methods for centering layer responses, gradients, and propagated errors, implemented by shortcut connections. In [44], an ``inception" layer is composed of a shortcut branch and a few deeper branches.

Concurrent with our work, ``highway networks" [42,43] present shortcut connections with gating functions [15]. These gates are data-dependent and have parameters, in contrast to our identity shortcuts that are parameter-free. When a gated shortcut is ``closed" (approaching zero), the layers in highway networks represent \emph{non-residual} functions. On the contrary, our formulation always learns residual functions; our identity shortcuts are never closed, and all information is always passed through, with additional residual functions to be learned. In addition, highway networks have not demonstrated accuracy gains with extremely increased depth (\eg, over 100 layers).



